\chapter{Writing Backend Services}
\label{ch:services}

Services are the heart of a Kiss backend.  A service is simply a class with methods --- no configuration files, no annotations, no registration.  This chapter covers how to write, organize, and debug services.

\section{The Service Pattern}
\label{sec:svc-pattern}

Every service method has the same four-parameter signature:

\begin{lstlisting}[style=groovy]
class MyService {
    void myMethod(JSONObject injson, JSONObject outjson,
                  Connection db, ProcessServlet servlet) {
        // Read input from injson
        // Do work (database queries, business logic)
        // Write output to outjson
    }
}
\end{lstlisting}

\noindent
The four parameters are:

\begin{description}[style=nextline]
  \item[\texttt{injson}] A \texttt{JSONObject} containing the data sent from the frontend.
  \item[\texttt{outjson}] A \texttt{JSONObject} that the method populates with the response.
  \item[\texttt{db}] A database \texttt{Connection} with auto-commit disabled.  The framework commits on success and rolls back on exception.
  \item[\texttt{servlet}] The \texttt{ProcessServlet} instance, used for file uploads and advanced features.
\end{description}

\noindent
The framework handles everything else: authentication, JSON parsing, database connection management, transaction control, and error responses.

\begin{tipbox}
A service method's only job is business logic.  It reads from \texttt{injson}, does work, and writes to \texttt{outjson}.
\end{tipbox}

\begin{notebox}
Although services are compiled on the fly when their source changes, they always run at full compiled bytecode speed --- never interpreted.  This applies equally to Groovy, Java, and Common Lisp services.  The compile step happens once per change; every subsequent request executes the cached bytecode at full JVM performance.  Hot reload is a development-time convenience with no runtime cost.
\end{notebox}

\section{Service Discovery}
\label{sec:svc-discovery}

Service files live in \texttt{src/main/backend/services/}.  The framework finds them by class name: a call to \texttt{Server.call("services/MyService",} \texttt{"myMethod", ...)} maps to \texttt{services/MyService.groovy} (or \texttt{.java} or \texttt{.lisp}).

The resolution order is:
\begin{enumerate}
  \item Groovy (\texttt{.groovy})
  \item Java (\texttt{.java})
  \item Common Lisp (\texttt{.lisp})
\end{enumerate}

\noindent
Subdirectories are supported.  A call to
\texttt{Server.call("services/\allowbreak{}admin/\allowbreak{}UserService", ...)} maps to
\texttt{services/admin/\allowbreak{}UserService.groovy}.

\section{Groovy Services}
\label{sec:svc-groovy}

Groovy is the most common language for Kiss services.  Files are automatically compiled on change --- no build step required.  The on-the-fly compile phase for Groovy is also faster than for Java, so the edit-save-reload cycle is noticeably quicker during development.

\begin{lstlisting}[style=groovy]
class Crud {
    void getRecords(JSONObject injson, JSONObject outjson,
                    Connection db, ProcessServlet servlet) {
        List<Record> recs = db.fetchAll(
            "select * from phone order by last_name, first_name")
        JSONArray rows = new JSONArray()
        for (Record rec : recs) {
            JSONObject row = new JSONObject()
            row.put("id", rec.getInt("rec_id"))
            row.put("firstName", rec.getString("first_name"))
            row.put("lastName", rec.getString("last_name"))
            row.put("phoneNumber", rec.getString(
                "phone_number"))
            rows.put(row)
        }
        outjson.put("rows", rows)
    }
}
\end{lstlisting}

\noindent
This service fetches all phone records, builds a JSON array, and returns it.  The frontend receives the \texttt{rows} array in the callback.

\begin{notebox}
Groovy's concise syntax means that semicolons are optional, \texttt{return} is optional for the last expression, and type declarations can use \texttt{def} instead of explicit types.  See Chapter~\ref{ch:groovy} for a full primer.
\end{notebox}

\section{Java Services}
\label{sec:svc-java}

Java services work identically but use \texttt{.java} files.  They are also auto-compiled on change during development.

\begin{lstlisting}[style=java]
import org.kissweb.json.JSONObject;
import org.kissweb.database.Connection;
import org.kissweb.restServer.ProcessServlet;

public class MyJavaService {
    public void addNumbers(JSONObject injson,
        JSONObject outjson,
        Connection db, ProcessServlet servlet) {
        Integer num1 = injson.getInt("num1");
        Integer num2 = injson.getInt("num2");
        if (num1 == null || num2 == null)
            outjson.put("num3", 0);
        else
            outjson.put("num3", num1 + num2);
    }
}
\end{lstlisting}

\section{Lisp Services}
\label{sec:svc-lisp}

Services can be written in Common Lisp using ABCL (Armed Bear Common Lisp).  This option is available for teams that prefer Lisp but is less commonly used.

\section{Reading Input and Writing Output}
\label{sec:svc-io}

Input arrives in the \texttt{injson} parameter.  Output is written to \texttt{outjson}.

\begin{lstlisting}[style=groovy]
// Reading input
String name = injson.getString("name")
int id = injson.getInt("id")
double amount = injson.getDouble("amount")
boolean active = injson.getBoolean("active")

// Writing output
outjson.put("status", "success")
outjson.put("count", 42)
outjson.put("total", 1234.56)
\end{lstlisting}

\noindent
Input values come directly from the JavaScript object passed to \texttt{Server.call()}.  Output values become properties of the response object in the JavaScript callback.

\subsection{Returning Record Lists}

The most common pattern is to return a list of database records as a JSON array:

\begin{lstlisting}[style=groovy]
// Option 1: Build JSON manually (full control over
// field names)
List<Record> recs = db.fetchAll("select * from employee")
JSONArray rows = new JSONArray()
for (Record rec : recs) {
    JSONObject row = new JSONObject()
    row.put("id", rec.getInt("employee_id"))
    row.put("name", rec.getString("last_name"))
    rows.put(row)
}
outjson.put("rows", rows)

// Option 2: Use fetchAllJSON (column names become JSON keys)
outjson.put("rows", db.fetchAllJSON(
    "select employee_id, last_name, first_name from employee"))
\end{lstlisting}

\begin{tipbox}
Use \texttt{fetchAllJSON()} when the database column names match what the frontend expects.  Use manual JSON building when you need to rename columns or compute derived values.
\end{tipbox}

\subsection{Returning a Single Record}

\begin{lstlisting}[style=groovy]
Record rec = db.fetchOne(
    "select * from employee where employee_id = ?", id)
if (rec == null)
    throw new UserException("Employee not found")
outjson.put("firstName", rec.getString("first_name"))
outjson.put("lastName", rec.getString("last_name"))
outjson.put("salary", rec.getDouble("salary"))
\end{lstlisting}

\section{Error Handling}
\label{sec:svc-errors}

Kiss provides three exception types for different error scenarios:

\begin{description}[style=nextline]
  \item[\texttt{UserException}]\index{UserException} Sends a message to the frontend for the user to see.  Does \emph{not} log a backend stack trace.  Use this for validation errors and expected conditions.

\begin{lstlisting}[style=groovy]
if (name == null || name.isEmpty())
    throw new UserException("Name is required")
\end{lstlisting}

  \item[\texttt{ServerException}]\index{ServerException} Logs the full stack trace in the backend log and sends a generic error code to the frontend.  Use this for unexpected errors where the user should not see implementation details.

  \item[\texttt{LogException}]\index{LogException} Logs the full stack trace \emph{and} sends the error message to the frontend.  Use this when both debugging information and a user-visible message are needed.
\end{description}

\noindent
If a service method throws any other exception, the framework catches it, logs the stack trace, rolls back the transaction, and sends a generic error response.

\section{A Complete CRUD Service}
\label{sec:svc-crud-example}

The following service demonstrates all four CRUD operations, a PDF report, and a CSV export:

\begin{lstlisting}[style=groovy]
package services

import org.kissweb.json.*
import org.kissweb.database.*
import org.kissweb.*
import org.kissweb.restServer.ProcessServlet

class Crud {

    void getRecords(JSONObject injson, JSONObject outjson,
                    Connection db, ProcessServlet servlet) {
        List<Record> recs = db.fetchAll(
            "select * from phone order by last_name, first_name")
        JSONArray rows = new JSONArray()
        for (Record rec : recs) {
            JSONObject row = new JSONObject()
            row.put("id", rec.getInt("rec_id"))
            row.put("firstName", rec.getString("first_name"))
            row.put("lastName", rec.getString("last_name"))
            row.put("phoneNumber", rec.getString(
                "phone_number"))
            rows.put(row)
        }
        outjson.put("rows", rows)
    }

    void addRecord(JSONObject injson, JSONObject outjson,
                   Connection db, ProcessServlet servlet) {
        Record rec = db.newRecord("phone")
        rec.set("first_name", injson.getString("firstName"))
        rec.set("last_name", injson.getString("lastName"))
        rec.set("phone_number", injson.getString(
            "phoneNumber"))
        rec.addRecord()
    }

    void updateRecord(JSONObject injson, JSONObject outjson,
                      Connection db, ProcessServlet servlet) {
        Record rec = db.fetchOne(
            "select * from phone where rec_id=?",
            injson.getInt("id"))
        rec.set("first_name", injson.getString("firstName"))
        rec.set("last_name", injson.getString("lastName"))
        rec.set("phone_number", injson.getString(
            "phoneNumber"))
        rec.update()
    }

    void deleteRecord(JSONObject injson, JSONObject outjson,
                      Connection db, ProcessServlet servlet) {
        db.execute("delete from phone where rec_id=?",
                   injson.getInt("id"))
    }

    void runReport(JSONObject injson, JSONObject outjson,
                   Connection db, ProcessServlet servlet) {
        Groff rpt = new Groff("PhoneList-", "Phone List",
            false)
        rpt.startTable("Lw(1i) Lw(1i) Lw(1i)")
        rpt.column("Last Name")
        rpt.column("First Name")
        rpt.column("Phone Number")
        rpt.endTitle()

        List<Record> recs = db.fetchAll(
            "select * from phone order by last_name, first_name")
        for (Record rec : recs) {
            rpt.column(rec.getString("last_name"))
            rpt.column(rec.getString("first_name"))
            rpt.column(rec.getString("phone_number"))
        }
        outjson.put("reportUrl", rpt.process())
    }

    void runExport(JSONObject injson, JSONObject outjson,
                   Connection db, ProcessServlet servlet) {
        File f = FileUtils.createReportFile("PhoneList-",
            ".csv")
        DelimitedFileWriter dfw =
            new DelimitedFileWriter(f.getAbsolutePath())

        dfw.writeField("Last Name")
        dfw.writeField("First Name")
        dfw.writeField("Phone Number")
        dfw.endRecord()

        List<Record> recs = db.fetchAll(
            "select * from phone order by last_name, first_name")
        for (Record rec : recs) {
            dfw.writeField(rec.getString("last_name"))
            dfw.writeField(rec.getString("first_name"))
            dfw.writeField(rec.getString("phone_number"))
            dfw.endRecord()
        }
        dfw.close()
        outjson.put("exportUrl", FileUtils.getHTTPPath(f))
    }
}
\end{lstlisting}

\begin{notebox}
This is the actual \texttt{Crud.groovy} from the Kiss example application.  Notice:
\begin{itemize}
  \item No configuration files or annotations.
  \item No explicit transaction management (the framework handles it).
  \item No authentication code (the framework checks the session automatically).
  \item Six methods, each a complete JSON-RPC endpoint.
\end{itemize}
\end{notebox}

\section{File Upload Handling}
\label{sec:svc-upload}

File uploads use the \texttt{servlet} parameter to access uploaded files:

\begin{lstlisting}[style=groovy]
void upload(JSONObject injson, JSONObject outjson,
            Connection db, ProcessServlet servlet) {
    int numberOfFiles = servlet.getUploadFileCount()
    String fileName = servlet.getUploadFileName(0)
    BufferedInputStream is
        = servlet.getUploadBufferedInputStream(0)
    // Process the input stream (save to disk, parse, etc.)
    is.close()
}
\end{lstlisting}

\noindent
See Chapter~\ref{ch:files} for complete file upload examples.

\section{Interfacing with External Systems}
\label{sec:svc-external-pointer}

Services frequently need to call out to third-party APIs or accept inbound calls from external systems that don't speak Kiss's JSON-RPC.  Both directions, plus exposing the application as an MCP server to AI assistants, are covered in Chapter~\ref{ch:external}.

\section{Calling Other Dynamically Compiled Services}
\label{sec:svc-cross-call}

Dynamically compiled services (Groovy, Java, or Common Lisp) can freely call any precompiled code in the system --- framework classes such as \texttt{Connection}, \texttt{Record}, \texttt{RestClient}, the Kiss utility classes, and anything you place in \texttt{src/main/precompiled/} --- using ordinary imports and method calls.  Within a single source file, methods, inner classes, and helpers can also call one another directly with no special handling.

What you \emph{cannot} do is have one dynamically compiled service hold a direct, compile-time reference to a class defined in a \emph{different} dynamically compiled source file.  Each on-the-fly compile produces a fresh class identity, so any direct reference would become stale the moment the target file is edited and recompiled --- defeating hot reload.  The framework therefore requires cross-file calls to go through an indirection that always resolves to the most recently compiled version of the target.

For Groovy services, that indirection is \texttt{GroovyService.run} (with \texttt{GroovyService.getMethod} as a companion for repeated calls):

\begin{lstlisting}[style=groovy]
import org.kissweb.restServer.GroovyService

// Single call: latest version is loaded, compiled if needed,
// and the method is invoked.
Object result = GroovyService.run(
    "services",        // path relative to src/main/backend/
    "OtherService",    // class/file name (no .groovy)
    "doWork",          // method name
    null,              // instance, or null for static
    arg1, arg2)        // boxed or unboxed arguments
\end{lstlisting}

\noindent
When the same method will be called repeatedly, fetch the \texttt{Method} handle once and reuse it:

\begin{lstlisting}[style=groovy]
import java.lang.reflect.Method
import org.kissweb.restServer.GroovyService

Method m = GroovyService.getMethod(
    "services", "OtherService", "doWork",
    String.class, Integer.class)

Object r1 = GroovyService.run(m, null, "alpha", 1)
Object r2 = GroovyService.run(m, null, "beta",  2)
\end{lstlisting}

\begin{notebox}
\texttt{GroovyService.run} verifies that the target source file is loaded and up-to-date before dispatching the call, so hot reload continues to work seamlessly across service boundaries.  Direct references would not.
\end{notebox}

\begin{tipbox}
This restriction is \emph{only} for crossing source-file boundaries between dynamically compiled units.  Methods within the same source file call each other normally, and any dynamically compiled service can call precompiled (framework or \texttt{precompiled/}) code directly.
\end{tipbox}

\section{Reusable Scripts}
\label{sec:svc-scripts}

The same cross-file mechanism is the standard way to share logic between services via reusable scripts.  Groovy scripts in \texttt{src/main/backend/scripts/} can be called from any service:

\begin{lstlisting}[style=groovy]
// In scripts/MyScript.groovy
public static Object myMethod(Integer a, Integer b) {
    return a + b + 100
}
\end{lstlisting}

\begin{lstlisting}[style=groovy]
// In a service method
import org.kissweb.restServer.GroovyService

Object result = GroovyService.run(false, "~/scripts",
    "MyScript", "myMethod", null, num1, num2)
\end{lstlisting}

\section{Services Without a Database}
\label{sec:svc-nodb}

If the application has no database configured, the \texttt{db} parameter will be \texttt{null}.  Services should check for this:

\begin{lstlisting}[style=groovy]
void myMethod(JSONObject injson, JSONObject outjson,
              Connection db, ProcessServlet servlet) {
    if (db == null) {
        outjson.put("nodb", true)
        return
    }
    // ... normal database operations ...
}
\end{lstlisting}

\section{Long-Running Requests}
\label{sec:long-running-requests}

Most service methods return quickly.  A few --- a slow LLM generation, a large
report, a bulk export --- can run for many seconds or minutes.  Kiss processes
every request asynchronously, and the servlet container imposes a default
timeout (roughly thirty seconds with the bundled Tomcat) on each asynchronous
request.  When a method exceeds it, the container aborts the request in flight:
the front-end's call fails, and the user sees ``Error communicating with the
server.''

A method that may legitimately run long should lift that limit by calling
\texttt{setTimeout} on its \texttt{ProcessServlet} argument before the slow work
begins:

\begin{lstlisting}[style=groovy]
void bigReport(JSONObject injson, JSONObject outjson,
               Connection db, ProcessServlet servlet) {
    // no async timeout for this request
    servlet.setTimeout(0)
    // ... work that may take minutes ...
}
\end{lstlisting}

\noindent
A value of zero (or any value $\le 0$) means no timeout at all; a positive value
sets the timeout, in milliseconds, for this request only.  Other requests are
unaffected and keep the default as a safety net, so one stuck handler cannot
quietly tie up a worker forever.

\begin{notebox}
Lifting the servlet timeout governs only how long Kiss itself will wait for the
method to finish.  Work that calls out to something else --- a long database
query, an external HTTP request, an LLM server --- has its \emph{own} timeout
that must also be generous enough.  For the Ollama client specifically, see
``Timeouts and Long Generations'' in the LLM Integration chapter.
\end{notebox}

\section{Service Best Practices}
\label{sec:svc-bestpractices}

\begin{itemize}
  \item Keep services focused.  One service class per entity or functional area.
  \item Use \texttt{UserException} for validation errors the user should see.
  \item Use parameterized queries (\texttt{?} placeholders) for all SQL to prevent injection.
  \item Let the framework handle transactions.  Only call \texttt{db.commit()} explicitly for long-running batch operations.
  \item Do not set database columns that have defaults (e.g., \texttt{CURRENT\_TIMESTAMP}).
  \item Use Kiss utility classes (\texttt{DateTime}, \texttt{NumberUtils}, \texttt{FileUtils}) instead of standard Java/Groovy alternatives for consistent behavior.
\end{itemize}
